\section{Evaluation}
\label{sec:evaluation}

\subsection{Baselines}
\label{sec:baselines}

We evaluate FieldFormer against a broad set of baselines spanning probabilistic interpolation, coordinate-based neural fields, physics-informed neural fields, and recent mesh-dependent spatio-temporal reconstruction models. This evaluation is designed to compare FieldFormer both against methods that naturally operate on continuous coordinates and against strong grid- or sensor-topology-based models that must be adapted to sparse sensing regimes.

\noindent\textbf{Mesh-free baselines.}
We first compare against mesh-free methods that, like FieldFormer, can be queried at arbitrary spatio-temporal coordinates. Gaussian Processes, also referred to as Kriging in spatial statistics, provide a classical probabilistic interpolation baseline grounded in spatio-temporal kernels~\citep{cressie1990origins}. Since exact Gaussian process inference scales cubically with the number of observations, we use stochastic variational Gaussian processes (SVGPs)~\citep{hensman2013gaussian}, which approximate the posterior using inducing points while retaining flexible uncertainty-aware interpolation. In our RMSE/MAE tables, SVGP is evaluated using its posterior predictive mean, so these results measure point-prediction accuracy rather than the full quality of the predictive distribution.

We also include two coordinate-based neural field baselines. Fourier-MLPs~\citep{tancik2020fourier} use sinusoidal Fourier features of the input coordinates to represent high-frequency variation, while SIRENs~\citep{sitzmann2020implicit} use sinusoidal activations throughout the network to model oscillatory and fine-scale spatio-temporal structure. These methods learn continuous functions of the form $f(x,y,t)$ and can therefore be evaluated at arbitrary query locations, making them direct mesh-free comparators. We additionally evaluate a five-member Fourier-MLP deep ensemble trained with independent random seeds. The ensemble averages member predictions and is evaluated as a point predictor under the same RMSE/MAE protocol; although deep ensembles are often used for uncertainty-aware modeling, our tables report only ensemble-mean point accuracy.

\noindent\textbf{Physics-informed neural field baselines.}
When the governing physics is available, we additionally evaluate physics-regularized versions of the neural field baselines. Specifically, we train Fourier-MLP-PINN and SIREN-PINN variants by augmenting their data-fitting objectives with PINN-style residual losses~\citep{raissi2019physics}. These baselines test whether global coordinate neural fields benefit from explicit physical regularization when the PDE residual can be computed. They also provide an important comparison to FieldFormer's local physics-aware formulation: while Fourier-MLP-PINN and SIREN-PINN enforce physics through global coordinate functions, FieldFormer combines local neighborhood aggregation with optional PDE-based regularization.

\noindent\textbf{Mesh-dependent spatio-temporal baselines.}
Although FieldFormer is mesh-free, many recent reconstruction and imputation methods assume a fixed grid, rasterized field, or fixed sensor graph. To provide a stronger empirical comparison, we adapt three representative mesh-dependent baselines to our sparse sensing setting: RecFNO~\citep{zhao2024recfno}, ImputeFormer~\citep{nie2024imputeformer}, and Senseiver~\citep{santos2023development}.

RecFNO is adapted as a grid-based field reconstruction baseline. Sparse sensor values are rasterized onto the simulation grid at the nearest grid locations, and a corresponding binary mask is provided to indicate which entries are observed. The model input consists of the sparse value grid, the observation mask, and coordinate channels, and the FNO predicts a full grid-valued field. In our implementation, this is instantiated as a Voronoi-style FNO input representation, where sparse observations and masks are placed on the grid and processed by spectral convolution layers to produce full-field predictions.

ImputeFormer is adapted as a fixed-node masked-imputation baseline over the deployed sensor network. Since the original model is designed for structured spatio-temporal imputation rather than continuous coordinate querying, we represent the persistent sensor network as a fixed set of nodes and train the model over temporal windows. Inputs consist of observed values together with masks, and training randomly masks subsets of available entries so that the model learns to impute missing measurements over the fixed sensor-time array. This adaptation evaluates whether a transformer designed for fixed-node temporal imputation can recover missing sensor readings in the longitudinal sparse sensing regime.

Senseiver is adapted as a sensor-set encoder and query-decoder baseline. At each time step, available sensor measurements are encoded as a set of sensor tokens containing normalized coordinates, time, observed values, and observation masks. A latent bottleneck attends to the sensor set, and query tokens are decoded from the learned latent representation. Unlike FieldFormer, Senseiver does not construct query-specific local neighborhoods; instead, it relies on a globally regularized latent representation of the observed sensor set. This makes it a useful comparator for testing when global latent priors are beneficial relative to locality-aware reconstruction.

\noindent\textbf{Comparison scope.}
These mesh-dependent baselines require adaptations that are not needed by FieldFormer. RecFNO requires rasterizing sparse sensors onto a fixed grid, ImputeFormer assumes a fixed set of sensor nodes and temporal windows, and Senseiver compresses the observed sensor set into a latent representation before decoding queries. In contrast, FieldFormer operates directly on sparse coordinate-value observations and constructs adaptive local neighborhoods at query time. This mesh-free formulation gives FieldFormer several practical advantages: it supports arbitrary query coordinates, multi-resolution inference, irregular sensor layouts, and evolving sensor networks without redefining a grid, rebuilding a graph, or changing the model architecture. At the same time, including mesh-dependent baselines allows us to evaluate whether strong global or fixed-topology priors can compensate for sparse observations in regimes where global field recovery is underconstrained.


\subsection{Experimental Setup}
\label{sec:experimental_setup}

We evaluate all methods under sparse longitudinal sensing regimes, where observations are available at a limited number of spatial locations but over dense temporal traces. Our evaluation distinguishes between two reconstruction objectives: \emph{sensor-space imputation} and \emph{global field reconstruction}.

\noindent\textbf{Evaluation Regimes.}
In \emph{sensor-space imputation}, the deployed sensor topology is fixed and the goal is to predict held-out sensor-time observations from the same persistent sensor network. This setting evaluates whether a method can reconstruct missing or withheld measurements over the observational support induced by the deployed sensors. For FieldFormer, we ensure that held-out observations are not included in the local neighborhood context during training, so performance reflects genuine imputation rather than reuse of the target observation.

In \emph{sensor-holdout evaluation}, entire sensor locations are removed from the training set and used only for validation or testing. This setting probes whether a method can generalize to unseen spatial support. Since the target locations are never observed during training, this evaluation is substantially more underconstrained than sensor-space imputation and depends strongly on the inductive priors of each model. For real-world datasets where dense full-field ground truth is unavailable, sensor-holdout evaluation serves as a practical proxy for global spatial reconstruction, since it measures prediction accuracy away from the observed training sensors.

When dense simulation fields or gridded reference data are available, we additionally evaluate \emph{global field reconstruction} over the full spatial domain. This setting directly measures how each method extrapolates away from the observed sensor support. Because global reconstruction under extreme sparsity is generally underdetermined, we interpret these results as evaluating prior-guided surrogate reconstruction rather than identifiable recovery of the true latent field.

\noindent\textbf{Train/Validation/Test Splits.}
For each dataset, we construct train, validation, and test splits in ratio $0.8:0.1:0.1$ over the sparse observations. In the standard sensor-space setting, the split is performed over observed sensor-time pairs while preserving the deployed sensor network. Thus, all sensors may appear during training, but individual observations are held out for validation and testing. In the sensor-holdout setting, validation and test sensors are removed entirely from the training set, so evaluation is performed at spatial locations unseen during training. Validation performance is used for checkpoint selection and hyperparameter tuning, while all reported numbers are computed on the held-out test split.

\noindent\textbf{Input Representation and Normalization.}
All methods are trained using the same observed values, sensor masks, and data splits. For multivariate datasets, each output channel is normalized using statistics computed from the training observations only, and predictions are transformed back to the original scale before evaluation. Missing entries are excluded from both the training loss and metric computation using the corresponding observation masks. For atmospheric wind fields, we represent wind using Cartesian components rather than angular direction, avoiding discontinuities associated with circular variables.

\noindent\textbf{Training Protocol.}
We use a consistent training protocol across methods wherever possible. Models are trained with AdamW, gradient clipping, validation-based checkpoint selection, and early stopping. Hyperparameters are selected using validation performance under the same data split. Mesh-free baselines operate directly on coordinate-value observations. Mesh-dependent baselines are adapted to sparse sensing through rasterization, fixed-node temporal windows, or sensor-set conditioning, as described in Section~\ref{sec:baselines}. This allows us to compare FieldFormer against both coordinate-based models and recent grid- or topology-dependent spatio-temporal reconstruction methods.

\noindent\textbf{Metrics and Uncertainty.}
We report root mean squared error (RMSE) and mean absolute error (MAE) for both sensor-space prediction and global field reconstruction. Sensor-space RMSE and MAE are computed on held-out sparse sensor-time observations and serve as the primary metrics for persistent sparse sensor networks, where reliable prediction is concentrated around observational support. For datasets with dense reference fields, global RMSE and MAE are computed over the full spatial domain. For real-world datasets where dense ground-truth fields are unavailable, sensor-holdout evaluation serves as a proxy for spatial generalization beyond the observed training support.

For sensor-space metrics, we estimate uncertainty using bootstrap resampling with 1000 bootstrap samples and report mean $\pm$ standard deviation. Bootstrapping is performed over the sparse evaluation set to quantify variability in the estimated RMSE and MAE. These standard deviations measure uncertainty in the aggregate RMSE/MAE estimates under test-set resampling, not per-query predictive uncertainty, epistemic uncertainty, or uncertainty over plausible global field completions. For full-field reconstruction, we report RMSE and MAE over the available dense reference field, using subsampling when necessary for memory efficiency. 


\subsection{Datasets}
\label{sec:datasets}

\subsubsection{Synthetic PDE Benchmarks}
\label{sec:synthetic_datasets}

\noindent\textbf{Periodic Heat Equation.}
We first evaluate on a scalar diffusion-dominated benchmark governed by the anisotropic two-dimensional heat equation with periodic boundary conditions:
\[
u_t = \alpha_x u_{xx} + \alpha_y u_{yy} + f(x,y,t).
\]
The forcing term $f(x,y,t)$ is a smooth sinusoidal function of space and time, inducing anisotropic and spatially coupled dynamics. The domain is discretized on a $64 \times 64$ spatial grid with $10{,}000$ time steps, and numerical stability is enforced through the CFL condition. A smooth sinusoidal initial condition is evolved using explicit finite differences with periodic wrapping. From the resulting full field, we sample 20 sensor locations uniformly at random and extract noisy longitudinal sensor traces. This benchmark provides a controlled setting in which the governing equation, physical parameters, and boundary conditions are fully specified, allowing us to evaluate reconstruction under sparse sensing for a diffusion-like process with broad spatial coupling.

\noindent\textbf{Periodic Shallow Water Equations.}
We next evaluate on a vector-valued propagation-dominated benchmark governed by the linearized two-dimensional shallow water equations with periodic boundary conditions:
\[
\eta_t + H(u_x + v_y) = 0, \qquad 
u_t + g \eta_x = 0, \qquad 
v_t + g \eta_y = 0.
\]
These equations describe free-surface gravity waves with characteristic speed $c=\sqrt{gH}$. We initialize the system with a Gaussian bump in surface height, which generates outward-propagating wave dynamics over the domain. The fields $(\eta,u,v)$ are simulated on a $64 \times 64$ spatial grid for $10{,}000$ time steps using a numerically stable forward--backward scheme satisfying the CFL condition. We sample 20 sparse sensor locations and extract noisy time series from the simulated fields. This benchmark tests reconstruction of coupled multi-variable dynamics with oscillatory propagation and local wave-like dependencies.

\noindent\textbf{Advection--Diffusion Pollution Simulation.}
As a synthetic-but-realistic transport benchmark, we simulate pollutant dispersion over New Delhi using a two-dimensional advection--diffusion equation:
\[
u_t + \mathbf{v}(t)\cdot\nabla u
=
\kappa \nabla^2 u + S(x,y),
\]
on a normalized $[0,1]\times[0,1]$ spatial domain corresponding to the city region. The initial concentration field is constructed from measurements at 32 regulatory monitoring stations on June 1, 2018, interpolated to the simulation grid using Universal Kriging. The source field combines spatial layers representing industrial activity, brick kilns, population density, and traffic intensity. Advection is driven by a synthetic monsoon-mode wind field oriented toward the northeast, with diurnal variation, directional wobble, and autoregressive stochastic gusts. Diffusion is set to produce an advection-dominated regime at the city scale, and open boundaries are handled using radiation-style boundary conditions with a sponge layer. The PDE is discretized using upwind advection, central diffusion, and a two-stage Heun/RK2 time integrator. This benchmark captures heterogeneous sources, time-varying transport, and sparse noisy sensing, while still providing dense simulated reference fields for global reconstruction evaluation.

\subsubsection{Real-World Sparse Sensor Benchmarks}
\label{sec:real_datasets}

\noindent\textbf{Real-World Pollution Monitoring.}
We evaluate on a real-world air-quality monitoring dataset from New Delhi. The dataset uses the city's deployed monitoring network, consisting of 32 air-quality stations distributed over an area of approximately $858~\mathrm{km}^2$. These stations are operated by public monitoring agencies, including the Central Pollution Control Board (CPCB), Delhi Pollution Control Committee (DPCC), and the Indian Meteorological Department (IMD). We use hourly measurements collected over a 30-month period from May 1, 2018 to November 1, 2020. The prediction targets are PM$_{2.5}$ and PM$_{10}$ concentrations. This dataset represents a persistent sparse sensor network: measurements are temporally dense, but spatial coverage is limited to a small number of fixed monitoring locations. Since dense ground-truth concentration fields are unavailable, we evaluate sensor-space imputation and sensor-holdout generalization.

\noindent\textbf{Real-World Atmospheric Measurements.}
We also evaluate on a real-world atmospheric sensing dataset using the same New Delhi sensor locations and time period as the pollution monitoring dataset. Instead of pollutant concentrations, this dataset contains meteorological variables measured over the deployed monitoring network. We model four atmospheric variables: average temperature, relative humidity, and horizontal wind components derived from wind speed and direction. Representing wind as Cartesian components avoids discontinuities associated with angular direction. This benchmark tests multivariate sparse sensor modeling for coupled atmospheric fields, where variables exhibit different smoothness, transport, and temporal variation patterns. As with the real-world pollution dataset, dense full-field ground truth is unavailable, so we evaluate sensor-space imputation and use sensor-holdout evaluation as a proxy for spatial generalization beyond the observed training sensors.


\subsection{Results}
\label{sec:results}

\subsubsection{Sensor-Space Imputation}
\label{sec:sensor_space_results}

We first evaluate sensor-space imputation, where the goal is to reconstruct held-out sensor-time observations over a persistent sensing topology. This is the most observationally grounded evaluation regime under extreme sparsity, since predictions are restricted to locations supported by the deployed sensor network.

Table~\ref{tab:sensor_space_synthetic} reports sensor-space performance on the three synthetic PDE benchmarks. FieldFormer is consistently competitive across all three systems, achieving near-best performance on Heat, Pollution, and SWE. On Heat and Pollution, Fourier-MLP obtains slightly lower sparse RMSE/MAE, while FieldFormer remains close. On SWE, FieldFormer is nearly tied with SIREN on sparse sensor-space metrics while substantially outperforming several mesh-dependent and globally regularized baselines. These results indicate that locality-aware reconstruction is well suited to sparse sensor imputation across diffusion, wave, and transport regimes.

% \begin{table}[htbp]
% \centering
% \caption{Sensor-space imputation results on synthetic PDE benchmarks. Results are reported as mean $\pm$ std over bootstrap samples. Lower is better.}
% \label{tab:sensor_space_synthetic}
% \resizebox{0.8\textwidth}{!}{%
% \begin{tabular}{llcc}
% \hline
% Dataset & Model & RMSE & MAE \\
% \hline

% \multirow{9}{*}{Heat}
% & FieldFormer & 0.09888 $\pm$ 0.0007416 & 0.07876 $\pm$ 0.0005815 \\
% & Fourier-MLP & \textbf{0.09786 $\pm$ 0.0006985} & \textbf{0.07786 $\pm$ 0.0005517} \\
% & Fourier-MLP-PINN & 0.1584 $\pm$ 0.001396 & 0.1206 $\pm$ 0.0008432 \\
% & SIREN & 0.09883 $\pm$ 0.0006853 & 0.07877 $\pm$ 0.0004833 \\
% & SIREN-PINN & 0.1793 $\pm$ 0.0009729 & 0.1434 $\pm$ 0.0007186 \\
% & SVGP & 0.1185 $\pm$ 0.0008308 & 0.09252 $\pm$ 0.0006093 \\
% & RecFNO & 0.1038 $\pm$ 0.0006049 & 0.08269 $\pm$ 0.0004322 \\
% & ImputeFormer & 0.1058 $\pm$ 0.0006532 & 0.08423 $\pm$ 0.000568 \\
% & Senseiver & 0.1011 $\pm$ 0.0007449 & 0.08045 $\pm$ 0.0005726 \\
% \hline

% \multirow{9}{*}{Pollution}
% & FieldFormer & 0.1154 $\pm$ 0.00032 & 0.09216 $\pm$ 0.0002176 \\
% & Fourier-MLP & \textbf{0.1147 $\pm$ 0.0003252} & \textbf{0.09153 $\pm$ 0.0002235} \\
% & Fourier-MLP-PINN & 0.1148 $\pm$ 0.0003401 & 0.09159 $\pm$ 0.0002141 \\
% & SIREN & 0.1162 $\pm$ 0.0003095 & 0.0928 $\pm$ 0.0002089 \\
% & SIREN-PINN & 0.1158 $\pm$ 0.0002912 & 0.09245 $\pm$ 0.000185 \\
% & SVGP & 0.1171 $\pm$ 0.0003127 & 0.0934 $\pm$ 0.0003083 \\
% & RecFNO & 0.1173 $\pm$ 0.0002567 & 0.09375 $\pm$ 0.0002272 \\
% & ImputeFormer & 0.1185 $\pm$ 0.0003339 & 0.09444 $\pm$ 0.0002698 \\
% & Senseiver & 0.1163 $\pm$ 0.0002625 & 0.09294 $\pm$ 0.0001972 \\
% \hline

% \multirow{9}{*}{SWE}
% & FieldFormer & 0.04644 $\pm$ 0.000254 & 0.03703 $\pm$ 0.0001957 \\
% & Fourier-MLP & 0.1011 $\pm$ 0.0004408 & 0.08049 $\pm$ 0.0002947 \\
% & Fourier-MLP-PINN & 0.101 $\pm$ 0.0004342 & 0.08044 $\pm$ 0.0002824 \\
% & SIREN & \textbf{0.04623 $\pm$ 0.0002687} & \textbf{0.03683 $\pm$ 0.0001931} \\
% & SIREN-PINN & 0.1008 $\pm$ 0.0004355 & 0.08029 $\pm$ 0.000281 \\
% & SVGP & 0.101 $\pm$ 0.0004366 & 0.08044 $\pm$ 0.0002813 \\
% & RecFNO & 0.06296 $\pm$ 0.0005098 & 0.04939 $\pm$ 0.0004467 \\
% & ImputeFormer & 0.06512 $\pm$ 0.0005653 & 0.05111 $\pm$ 0.0004587 \\
% & Senseiver & 0.07722 $\pm$ 0.0005243 & 0.06056 $\pm$ 0.0003741 \\

% \hline
% \end{tabular}%
% }
% \end{table}

\begin{table}[htbp]
\centering
\caption{Sensor-space imputation results on synthetic PDE benchmarks. Results are reported as mean $\pm$ std over bootstrap samples. Lower is better.}
\label{tab:sensor_space_synthetic}
\resizebox{\textwidth}{!}{%
\begin{tabular}{lcccccc}
\hline
Model 
& Heat RMSE & Heat MAE
& Pollution RMSE & Pollution MAE
& SWE RMSE & SWE MAE \\
\hline

FieldFormer 
& 0.09888 $\pm$ 0.0007416 
& 0.07876 $\pm$ 0.0005815
& 0.1154 $\pm$ 0.00032
& 0.09216 $\pm$ 0.0002176
& 0.04644 $\pm$ 0.000254
& 0.03703 $\pm$ 0.0001957 \\

Fourier-MLP 
& \textbf{0.09786 $\pm$ 0.0006985}
& \textbf{0.07786 $\pm$ 0.0005517}
& \textbf{0.1147 $\pm$ 0.0003252}
& \textbf{0.09153 $\pm$ 0.0002235}
& 0.1011 $\pm$ 0.0004408
& 0.08049 $\pm$ 0.0002947 \\

Fourier-MLP Ensemble
& 1.588 $\pm$ 0.007539
& 1.124 $\pm$ 0.003787
& 0.1289 $\pm$ 0.0005315
& 0.1004 $\pm$ 0.0002781
& 0.1010 $\pm$ 0.0004365
& 0.08043 $\pm$ 0.0002840 \\

Fourier-MLP-PINN 
& 0.1584 $\pm$ 0.001396
& 0.1206 $\pm$ 0.0008432
& 0.1148 $\pm$ 0.0003401
& 0.09159 $\pm$ 0.0002141
& 0.101 $\pm$ 0.0004342
& 0.08044 $\pm$ 0.0002824 \\

SIREN 
& 0.09883 $\pm$ 0.0006853
& 0.07877 $\pm$ 0.0004833
& 0.1162 $\pm$ 0.0003095
& 0.0928 $\pm$ 0.0002089
& \textbf{0.04623 $\pm$ 0.0002687}
& \textbf{0.03683 $\pm$ 0.0001931} \\

SIREN-PINN 
& 0.1793 $\pm$ 0.0009729
& 0.1434 $\pm$ 0.0007186
& 0.1158 $\pm$ 0.0002912
& 0.09245 $\pm$ 0.000185
& 0.1008 $\pm$ 0.0004355
& 0.08029 $\pm$ 0.000281 \\

SVGP 
& 0.1185 $\pm$ 0.0008308
& 0.09252 $\pm$ 0.0006093
& 0.1171 $\pm$ 0.0003127
& 0.0934 $\pm$ 0.0003083
& 0.101 $\pm$ 0.0004366
& 0.08044 $\pm$ 0.0002813 \\

RecFNO 
& 0.1038 $\pm$ 0.0006049
& 0.08269 $\pm$ 0.0004322
& 0.1173 $\pm$ 0.0002567
& 0.09375 $\pm$ 0.0002272
& 0.06296 $\pm$ 0.0005098
& 0.04939 $\pm$ 0.0004467 \\

ImputeFormer 
& 0.1058 $\pm$ 0.0006532
& 0.08423 $\pm$ 0.000568
& 0.1185 $\pm$ 0.0003339
& 0.09444 $\pm$ 0.0002698
& 0.06512 $\pm$ 0.0005653
& 0.05111 $\pm$ 0.0004587 \\

Senseiver 
& 0.1011 $\pm$ 0.0007449
& 0.08045 $\pm$ 0.0005726
& 0.1163 $\pm$ 0.0002625
& 0.09294 $\pm$ 0.0001972
& 0.07722 $\pm$ 0.0005243
& 0.06056 $\pm$ 0.0003741 \\

\hline
\end{tabular}%
}
\end{table}

Table~\ref{tab:sensor_space_real} reports sensor-space imputation results on the real-world atmospheric and pollution monitoring datasets. FieldFormer performs especially strongly in this regime. On atmospheric measurements, FieldFormer achieves the best or near-best performance across temperature, relative humidity, and wind components, with particularly large gains over Fourier-MLP, the Fourier-MLP ensemble, SIREN, SVGP, ImputeFormer, and Senseiver on temperature and relative humidity. The ensemble modestly improves over single Fourier-MLP on most real-world sensor-space variables, but remains well behind FieldFormer. On the real-world pollution dataset, FieldFormer obtains the lowest RMSE and MAE for both PM$_{10}$ and PM$_{2.5}$, showing that locality-aware modeling is effective for persistent urban monitoring networks.

\begin{table}[htbp]
\centering
\caption{Sensor-space imputation results on real-world sparse sensor benchmarks. Results are reported as mean $\pm$ std over bootstrap samples. Lower is better.}
\label{tab:sensor_space_real}
\resizebox{\textwidth}{!}{%
\begin{tabular}{llcccc}
\hline
Dataset & Model & Variable 1 RMSE & Variable 1 MAE & Variable 2 RMSE & Variable 2 MAE \\
\hline

\multirow{8}{*}{Atmospheric: AT/RH}
& FieldFormer & \textbf{0.9966 $\pm$ 0.02692} & \textbf{0.5198 $\pm$ 0.004008} & \textbf{4.333 $\pm$ 0.0544} & \textbf{2.171 $\pm$ 0.01099} \\
& Fourier-MLP & 3.899 $\pm$ 0.01267 & 3.157 $\pm$ 0.009728 & 14.94 $\pm$ 0.03613 & 11.88 $\pm$ 0.03015 \\
& Fourier-MLP Ensemble & 3.878 $\pm$ 0.01237 & 3.144 $\pm$ 0.009327 & 14.87 $\pm$ 0.03755 & 11.83 $\pm$ 0.03070 \\
& SIREN & 4.125 $\pm$ 0.01249 & 3.333 $\pm$ 0.009027 & 16.28 $\pm$ 0.03885 & 13.03 $\pm$ 0.03152 \\
& SVGP & 8.494 $\pm$ 0.02936 & 6.905 $\pm$ 0.02905 & 23.47 $\pm$ 0.06189 & 19.42 $\pm$ 0.05223 \\
& RecFNO & 1.449 $\pm$ 0.02407 & 0.9899 $\pm$ 0.005657 & 5.992 $\pm$ 0.05269 & 3.694 $\pm$ 0.015 \\
& ImputeFormer & 1.619 $\pm$ 0.02632 & 1.056 $\pm$ 0.005216 & 6.978 $\pm$ 0.06032 & 4.346 $\pm$ 0.02544 \\
& Senseiver & 2.243 $\pm$ 0.02496 & 1.557 $\pm$ 0.009949 & 9.56 $\pm$ 0.08187 & 6.341 $\pm$ 0.02758 \\
\hline

\multirow{8}{*}{Atmospheric: $U_x$/$V_y$}
& FieldFormer & 0.6807 $\pm$ 0.0134 & \textbf{0.3213 $\pm$ 0.002875} & \textbf{0.7011 $\pm$ 0.01045} & \textbf{0.3468 $\pm$ 0.002512} \\
& Fourier-MLP & 0.8519 $\pm$ 0.008105 & 0.5463 $\pm$ 0.002747 & 0.8375 $\pm$ 0.00963 & 0.5132 $\pm$ 0.002926 \\
& Fourier-MLP Ensemble & 0.8412 $\pm$ 0.007255 & 0.5357 $\pm$ 0.002651 & 0.8290 $\pm$ 0.009320 & 0.5049 $\pm$ 0.002972 \\
& SIREN & 1.041 $\pm$ 0.005132 & 0.702 $\pm$ 0.002925 & 0.959 $\pm$ 0.007708 & 0.5865 $\pm$ 0.00253 \\
& SVGP & 1.357 $\pm$ 0.01282 & 0.8861 $\pm$ 0.004546 & 1.086 $\pm$ 0.009182 & 0.6687 $\pm$ 0.00327 \\
& RecFNO & \textbf{0.6653 $\pm$ 0.01107} & 0.3558 $\pm$ 0.002152 & 0.7204 $\pm$ 0.007887 & 0.3801 $\pm$ 0.001811 \\
& ImputeFormer & 0.7271 $\pm$ 0.01181 & 0.388 $\pm$ 0.002308 & 0.7361 $\pm$ 0.01278 & 0.399 $\pm$ 0.002718 \\
& Senseiver & 0.8421 $\pm$ 0.01156 & 0.5182 $\pm$ 0.00315 & 0.8082 $\pm$ 0.009127 & 0.5003 $\pm$ 0.002084 \\
\hline

\multirow{8}{*}{Pollution: PM$_{10}$/PM$_{2.5}$}
& FieldFormer & \textbf{39.25 $\pm$ 0.8496} & \textbf{20.28 $\pm$ 0.1172} & \textbf{22.06 $\pm$ 0.664} & \textbf{11.23 $\pm$ 0.0691} \\
& Fourier-MLP & 83.48 $\pm$ 0.3978 & 54.81 $\pm$ 0.1662 & 53.19 $\pm$ 0.4248 & 31.74 $\pm$ 0.1232 \\
& Fourier-MLP Ensemble & 81.40 $\pm$ 0.3946 & 53.25 $\pm$ 0.1711 & 51.80 $\pm$ 0.4218 & 30.75 $\pm$ 0.1193 \\
& SIREN & 92.81 $\pm$ 0.3154 & 62.41 $\pm$ 0.1189 & 58.95 $\pm$ 0.4278 & 35.93 $\pm$ 0.1205 \\
& SVGP & 147.9 $\pm$ 0.5388 & 110 $\pm$ 0.3128 & 98.19 $\pm$ 0.4255 & 67.41 $\pm$ 0.2203 \\
& RecFNO & 53.84 $\pm$ 0.7037 & 31.25 $\pm$ 0.19 & 32.91 $\pm$ 0.7806 & 17.65 $\pm$ 0.1067 \\
& ImputeFormer & 58.96 $\pm$ 0.6289 & 35.2 $\pm$ 0.1462 & 36.27 $\pm$ 0.7988 & 19.09 $\pm$ 0.1309 \\
& Senseiver & 60.05 $\pm$ 0.5948 & 36.94 $\pm$ 0.1758 & 36.21 $\pm$ 0.6968 & 19.98 $\pm$ 0.08917 \\

\hline
\end{tabular}%
}
\end{table}

Overall, the sensor-space results support the main intended use case of FieldFormer: imputation over persistent sparse sensor networks. FieldFormer is not uniformly best on every synthetic sensor-space metric, but it is consistently competitive and performs particularly strongly on real-world sparse sensing datasets, where fixed sensor topology and local spatial structure are central.

\subsubsection{Global Field Reconstruction}
\label{sec:global_reconstruction_results}

We next evaluate global field reconstruction, where the goal is to predict values away from the observed sensor support. For synthetic datasets, dense reference fields are available, allowing direct full-field evaluation. Table~\ref{tab:global_synthetic} merges the full-field results for Heat, Pollution, and SWE.

The results show that global reconstruction is substantially more variable than sensor-space imputation. No single method dominates across all datasets and metrics. FieldFormer remains competitive across all three benchmarks, but different baselines are favored in different regimes: RecFNO performs well on Heat full-field reconstruction, Fourier-MLP and SVGP are competitive on Pollution, and Fourier-MLP-PINN/SVGP-style smooth reconstructions perform strongly on SWE full-field metrics. This variability supports our central argument that global reconstruction under extreme sparse sensing is prior-dependent: away from observed sensors, performance depends strongly on how each model's inductive bias aligns with the underlying data-generating process.

% \begin{table}[htbp]
% \centering
% \caption{Global field reconstruction results on synthetic PDE benchmarks. Dense reference fields are available for these datasets. Lower is better.}
% \label{tab:global_synthetic}
% \resizebox{0.8\textwidth}{!}{%
% \begin{tabular}{llcc}
% \hline
% Dataset & Model & Full RMSE & Full MAE \\
% \hline

% \multirow{9}{*}{Heat}
% & FieldFormer & 0.1818 $\pm$ 0.001179 & 0.1313 $\pm$ 0.0008449 \\
% & Fourier-MLP & 0.2146 $\pm$ 0.001386 & 0.1469 $\pm$ 0.0009123 \\
% & Fourier-MLP-PINN & 0.2141 $\pm$ 0.001346 & 0.1517 $\pm$ 0.001026 \\
% & SIREN & 0.1783 $\pm$ 0.0009598 & 0.1328 $\pm$ 0.0007279 \\
% & SIREN-PINN & 0.195 $\pm$ 0.0009061 & 0.1533 $\pm$ 0.0007207 \\
% & SVGP & 0.1788 $\pm$ 0.001214 & 0.1306 $\pm$ 0.0008949 \\
% & RecFNO & \textbf{0.1771 $\pm$ 0.001173} & \textbf{0.1212 $\pm$ 0.0007414} \\
% & ImputeFormer & 0.1916 $\pm$ 0.001001 & 0.1435 $\pm$ 0.000781 \\
% & Senseiver & 0.1959 $\pm$ 0.002121 & 0.1421 $\pm$ 0.001958 \\
% \hline

% \multirow{9}{*}{Pollution}
% & FieldFormer & 0.4359 $\pm$ 0.04525 & 0.0831 $\pm$ 0.006115 \\
% & Fourier-MLP & \textbf{0.4326 $\pm$ 0.04487} & 0.09857 $\pm$ 0.005979 \\
% & Fourier-MLP-PINN & 0.4652 $\pm$ 0.04322 & 0.1314 $\pm$ 0.006371 \\
% & SIREN & 0.4556 $\pm$ 0.04306 & 0.1256 $\pm$ 0.006149 \\
% & SIREN-PINN & 0.4598 $\pm$ 0.0429 & 0.122 $\pm$ 0.006211 \\
% & SVGP & 0.4349 $\pm$ 0.04505 & \textbf{0.08208 $\pm$ 0.006107} \\
% & RecFNO & 0.5676 $\pm$ 0.03491 & 0.1684 $\pm$ 0.006464 \\
% & ImputeFormer & 0.4482 $\pm$ 0.0439 & 0.1103 $\pm$ 0.00613 \\
% & Senseiver & 0.4429 $\pm$ 0.00875 & 0.09243 $\pm$ 0.001077 \\
% \hline

% \multirow{9}{*}{SWE}
% & FieldFormer & 0.2015 $\pm$ 0.0002546 & 0.1465 $\pm$ 0.0002115 \\
% & Fourier-MLP & 0.1807 $\pm$ 0.0001687 & 0.1319 $\pm$ 0.0001417 \\
% & Fourier-MLP-PINN & 0.18 $\pm$ 0.0001695 & 0.1311 $\pm$ 0.0001448 \\
% & SIREN & 0.6818 $\pm$ 0.000437 & 0.473 $\pm$ 0.0003094 \\
% & SIREN-PINN & \textbf{0.1799 $\pm$ 0.0001696} & \textbf{0.131 $\pm$ 0.0001448} \\
% & SVGP & 0.18 $\pm$ 0.0001695 & 0.1311 $\pm$ 0.0001449 \\
% & RecFNO & 0.1985 $\pm$ 0.0002363 & 0.15 $\pm$ 0.0001697 \\
% & ImputeFormer & 0.1816 $\pm$ 0.0001737 & 0.1321 $\pm$ 0.0001524 \\
% & Senseiver & 0.1867 $\pm$ 0.0006654 & 0.1364 $\pm$ 0.0004551 \\

% \hline
% \end{tabular}%
% }
% \end{table}

\begin{table}[htbp]
\centering
\caption{Global field reconstruction results on synthetic PDE benchmarks. Dense reference fields are available for these datasets. Lower is better.}
\label{tab:global_synthetic}
\resizebox{\textwidth}{!}{%
\begin{tabular}{lcccccc}
\hline
Model 
& Heat RMSE & Heat MAE
& Pollution RMSE & Pollution MAE
& SWE RMSE & SWE MAE \\
\hline

FieldFormer 
& 0.1818 $\pm$ 0.001179 
& 0.1313 $\pm$ 0.0008449
& 0.4359 $\pm$ 0.04525
& 0.0831 $\pm$ 0.006115
& 0.2015 $\pm$ 0.0002546
& 0.1465 $\pm$ 0.0002115 \\

Fourier-MLP 
& 0.2146 $\pm$ 0.001386
& 0.1469 $\pm$ 0.0009123
& \textbf{0.4326 $\pm$ 0.04487}
& 0.09857 $\pm$ 0.005979
& 0.1807 $\pm$ 0.0001687
& 0.1319 $\pm$ 0.0001417 \\

Fourier-MLP Ensemble
& 1.326 $\pm$ 0.008649
& 0.9717 $\pm$ 0.006706
& 0.4344 $\pm$ 0.04490
& 0.09691 $\pm$ 0.006017
& 0.1828 $\pm$ 0.0001669
& 0.1340 $\pm$ 0.0001324 \\

Fourier-MLP-PINN 
& 0.2141 $\pm$ 0.001346
& 0.1517 $\pm$ 0.001026
& 0.4652 $\pm$ 0.04322
& 0.1314 $\pm$ 0.006371
& 0.18 $\pm$ 0.0001695
& 0.1311 $\pm$ 0.0001448 \\

SIREN 
& 0.1783 $\pm$ 0.0009598
& 0.1328 $\pm$ 0.0007279
& 0.4556 $\pm$ 0.04306
& 0.1256 $\pm$ 0.006149
& 0.6818 $\pm$ 0.000437
& 0.473 $\pm$ 0.0003094 \\

SIREN-PINN 
& 0.195 $\pm$ 0.0009061
& 0.1533 $\pm$ 0.0007207
& 0.4598 $\pm$ 0.0429
& 0.122 $\pm$ 0.006211
& \textbf{0.1799 $\pm$ 0.0001696}
& \textbf{0.131 $\pm$ 0.0001448} \\

SVGP 
& 0.1788 $\pm$ 0.001214
& 0.1306 $\pm$ 0.0008949
& 0.4349 $\pm$ 0.04505
& \textbf{0.08208 $\pm$ 0.006107}
& 0.18 $\pm$ 0.0001695
& 0.1311 $\pm$ 0.0001449 \\

RecFNO 
& \textbf{0.1771 $\pm$ 0.001173}
& \textbf{0.1212 $\pm$ 0.0007414}
& 0.5676 $\pm$ 0.03491
& 0.1684 $\pm$ 0.006464
& 0.1985 $\pm$ 0.0002363
& 0.15 $\pm$ 0.0001697 \\

ImputeFormer 
& 0.1916 $\pm$ 0.001001
& 0.1435 $\pm$ 0.000781
& 0.4482 $\pm$ 0.0439
& 0.1103 $\pm$ 0.00613
& 0.1816 $\pm$ 0.0001737
& 0.1321 $\pm$ 0.0001524 \\

Senseiver 
& 0.1959 $\pm$ 0.002121
& 0.1421 $\pm$ 0.001958
& 0.4429 $\pm$ 0.00875
& 0.09243 $\pm$ 0.001077
& 0.1867 $\pm$ 0.0006654
& 0.1364 $\pm$ 0.0004551 \\

\hline
\end{tabular}%
}
\end{table}

For real-world datasets, dense full-field ground truth is unavailable. We therefore use sensor-holdout evaluation as a proxy for spatial generalization: entire sensors are removed from training and used only for testing. Table~\ref{tab:global_proxy_real} reports these results for atmospheric and pollution monitoring datasets. This setting should be interpreted with some care because the baselines expose different inductive biases and query interfaces. In particular, ImputeFormer is a fixed-node imputer over the deployed sensor array, giving it access to a transductive fixed-topology structure that is advantageous when the held-out targets are drawn from the same persistent sensor network.

The real-world sensor-holdout results further demonstrate the prior-dependent nature of spatial extrapolation. ImputeFormer performs strongly in this setting, especially on relative humidity, wind-$U_x$, and both pollution variables, but this should be read together with its fixed-node/transductive advantage. Senseiver is also competitive, obtaining the best atmospheric temperature RMSE, but it is not consistently best across variables or metrics. Ensembling the Fourier-MLP helps on some sensor-holdout metrics but does not produce a uniformly strong extrapolator. FieldFormer is less dominant in this setting than in sensor-space imputation, consistent with the fact that sensor-holdout evaluation removes the local observational support that locality-aware methods rely on. Overall, Table~\ref{tab:global_proxy_real} reinforces the same conclusion as the synthetic full-field results: no model uniformly solves reconstruction away from observed support; performance depends on how each model's prior matches the dataset and evaluation regime.

\begin{table}[htbp]
\centering
\caption{Sensor-holdout results on real-world sparse sensor benchmarks, used as a proxy for spatial generalization when dense full-field ground truth is unavailable. Results are reported as mean $\pm$ std over bootstrap samples. Lower is better. ImputeFormer is a fixed-node imputer, giving it a transductive fixed-topology advantage in this sensor-holdout setting.}
\label{tab:global_proxy_real}
\resizebox{\textwidth}{!}{%
\begin{tabular}{llcccc}
\hline
Dataset & Model & Variable 1 RMSE & Variable 1 MAE & Variable 2 RMSE & Variable 2 MAE \\
\hline

\multirow{8}{*}{Atmospheric: AT/RH}
& FieldFormer & 4.63 $\pm$ 0.6376 & 3.14 $\pm$ 0.4384 & 15.75 $\pm$ 2.517 & 10.04 $\pm$ 1.568 \\
& Fourier-MLP & 5.403 $\pm$ 0.3827 & 4.367 $\pm$ 0.2752 & 27.33 $\pm$ 2.096 & 22.55 $\pm$ 1.938 \\
& Fourier-MLP Ensemble & 5.091 $\pm$ 0.3700 & 4.039 $\pm$ 0.2682 & 18.68 $\pm$ 0.9145 & 15.31 $\pm$ 0.8338 \\
& SIREN & 6.353 $\pm$ 0.4423 & 5.045 $\pm$ 0.3822 & 21.16 $\pm$ 1.12 & 17.27 $\pm$ 1.018 \\
& SVGP & 8.795 $\pm$ 0.5025 & 7.144 $\pm$ 0.4788 & 22.01 $\pm$ 0.9019 & 18.54 $\pm$ 0.8679 \\
& RecFNO & 11.34 $\pm$ 0.7088 & 9.333 $\pm$ 0.59 & 21.88 $\pm$ 1.137 & 18.2 $\pm$ 1.039 \\
& ImputeFormer & 3.037 $\pm$ 0.6049 & \textbf{2.001 $\pm$ 0.3392} & \textbf{7.208 $\pm$ 0.611} & \textbf{5.327 $\pm$ 0.3565} \\
& Senseiver & \textbf{3.017 $\pm$ 0.5426} & 2.042 $\pm$ 0.3465 & 10.59 $\pm$ 1.065 & 8.112 $\pm$ 0.8772 \\
\hline

\multirow{8}{*}{Atmospheric: $U_x$/$V_y$}
& FieldFormer & 2.26 $\pm$ 0.6081 & 1.282 $\pm$ 0.2269 & 1.502 $\pm$ 0.1426 & 1.026 $\pm$ 0.09529 \\
& Fourier-MLP & 2.162 $\pm$ 0.7587 & 1.183 $\pm$ 0.2525 & 0.9047 $\pm$ 0.1635 & 0.5812 $\pm$ 0.0651 \\
& SIREN & 2.101 $\pm$ 0.7676 & 1.027 $\pm$ 0.2498 & 0.9368 $\pm$ 0.1635 & 0.614 $\pm$ 0.0696 \\
& Fourier-MLP Ensemble & 2.015 $\pm$ 0.7490 & 0.9600 $\pm$ 0.2368 & \textbf{0.8970 $\pm$ 0.1819} & 0.5455 $\pm$ 0.07301 \\
& SVGP & 2.033 $\pm$ 0.7474 & 0.9731 $\pm$ 0.2393 & 0.9001 $\pm$ 0.1931 & \textbf{0.5331 $\pm$ 0.07598} \\
& RecFNO & 2.034 $\pm$ 0.7 & 0.9904 $\pm$ 0.2369 & 1.014 $\pm$ 0.1774 & 0.6318 $\pm$ 0.07666 \\
& ImputeFormer & \textbf{1.911 $\pm$ 0.8183} & \textbf{0.7754 $\pm$ 0.2425} & 0.9473 $\pm$ 0.1837 & 0.5974 $\pm$ 0.08248 \\
& Senseiver & 2.091 $\pm$ 0.687 & 1.001 $\pm$ 0.2517 & 1.447 $\pm$ 0.2026 & 0.9016 $\pm$ 0.1555 \\
\hline

\multirow{8}{*}{Pollution: PM$_{10}$/PM$_{2.5}$}
& FieldFormer & 115 $\pm$ 13.11 & 69.63 $\pm$ 9.17 & 45.95 $\pm$ 4.77 & 26.76 $\pm$ 3.152 \\
& Fourier-MLP & 121.7 $\pm$ 9.278 & 82.31 $\pm$ 6.506 & 67.77 $\pm$ 6.211 & 42.41 $\pm$ 4.393 \\
& Fourier-MLP Ensemble & 120.8 $\pm$ 8.308 & 82.58 $\pm$ 6.140 & 68.39 $\pm$ 5.787 & 44.01 $\pm$ 4.061 \\
& SIREN & 149.4 $\pm$ 12.78 & 103.6 $\pm$ 8.571 & 82.73 $\pm$ 7.859 & 54.26 $\pm$ 4.981 \\
& SVGP & 166.5 $\pm$ 12.13 & 122 $\pm$ 7.467 & 99.17 $\pm$ 9.501 & 68.2 $\pm$ 6.101 \\
& RecFNO & 170.9 $\pm$ 15.06 & 127.4 $\pm$ 11.63 & 75.75 $\pm$ 5.656 & 52.74 $\pm$ 2.747 \\
& ImputeFormer & \textbf{80.76 $\pm$ 10.07} & \textbf{47.09 $\pm$ 6.202} & \textbf{38.93 $\pm$ 4.384} & \textbf{21.15 $\pm$ 2.274} \\
& Senseiver & 86.51 $\pm$ 11.83 & 50.6 $\pm$ 7.264 & 39.84 $\pm$ 4.632 & 22.77 $\pm$ 2.687 \\

\hline
\end{tabular}%
}
\end{table}

Taken together, the global reconstruction and sensor-holdout results show that extrapolating away from observed sensor support is substantially more sensitive to the choice of inductive prior than imputing over an existing sensor network. FieldFormer performs strongly when local sensor support is available, fixed-node imputers can benefit from transductive topology in sensor-holdout settings, and globally regularized baselines can be advantageous when their priors align with the dataset. The absence of a uniform winner supports our broader framing: under extreme sparse sensing, global reconstruction should be treated as prior-guided surrogate completion, whereas sensor-space imputation is the more identifiable and practically reliable objective.

\paragraph{Inference efficiency.}
Appendix~\ref{sec:inference_efficiency} reports wall-clock sparse-test prediction time per query, peak GPU memory, and setup time using the same trained checkpoints and sparse-test protocol. The efficiency comparison is most meaningful against context-conditioned reconstruction models, since direct coordinate fields such as Fourier-MLP and SIREN do not aggregate observed sensor context and therefore do not incur attention or set-conditioning costs. In this comparable regime, FieldFormer is an efficient transformer architecture: its fixed-size local attention is consistently over $100\times$ faster per query than Senseiver's global sensor-set transformer and roughly $7$--$10\times$ faster than the grid-based RecFNO evaluator across the timing benchmarks, while retaining mesh-free coordinate querying. ImputeFormer has lower query-time latency after setup, but this is because its fixed-node/transductive formulation precomputes and caches predictions over the deployed sensor-time array, reducing evaluation to lookup rather than arbitrary coordinate reconstruction. These results support FieldFormer's design goal: efficient transformer-based sparse reconstruction through local context rather than global attention over the observation set.

\subsection{Ablation Studies}
\label{sec:ablations}

We ablate three FieldFormer variants to isolate the role of physics regularization, transformer-based local aggregation, and spatially adaptive neighborhood geometry. \textbf{FieldFormer-PINN} augments FieldFormer with the PDE residual and boundary losses described in Section~\ref{sec:ff_variants}. \textbf{FieldFormer-Gamma Field} replaces global velocity-scaling parameters with coordinate-dependent gamma fields, also described in Section~\ref{sec:ff_variants}. \textbf{FieldFormer-MLP} preserves the same mesh-free neighborhood construction and velocity-scaled spatio-temporal metric as FieldFormer, but replaces the local transformer encoder with a multilayer perceptron. In this variant, the selected local neighborhood is encoded through relative offsets and observed values, then aggregated by an MLP prediction head rather than attention. This ablation tests whether performance gains arise primarily from adaptive local evidence selection or from transformer-based relational reasoning within the neighborhood.

\paragraph{Synthetic PDE benchmarks.}
Table~\ref{tab:ablation_synthetic_sparse} reports sensor-space imputation results on the three synthetic PDE benchmarks. FieldFormer without explicit physics regularization performs best on Heat and SWE sensor-space prediction, while FieldFormer-Gamma Field is marginally strongest on the Pollution simulation. FieldFormer-MLP is competitive on Heat but degrades more noticeably on Pollution and SWE, suggesting that attention-based aggregation is useful when local neighborhoods contain coupled or directional structure. FieldFormer-PINN generally worsens sensor-space performance, especially on Heat and SWE, indicating that physics residuals can conflict with sparse noisy observations when the reconstruction target is sensor-space imputation.

\begin{table}[htbp]
\centering
\caption{Ablation study on synthetic PDE benchmarks: sensor-space imputation. Results are reported as mean $\pm$ std over bootstrap samples. Lower is better.}
\label{tab:ablation_synthetic_sparse}
\resizebox{\textwidth}{!}{%
\begin{tabular}{lcccccc}
\hline
Model 
& Heat RMSE & Heat MAE
& Pollution RMSE & Pollution MAE
& SWE RMSE & SWE MAE \\
\hline

FieldFormer-PINN
& 0.1071 $\pm$ 0.0009743
& 0.08513 $\pm$ 0.000793
& 0.1214 $\pm$ 0.0003221
& 0.09683 $\pm$ 0.0002315
& 0.05716 $\pm$ 0.0002378
& 0.04541 $\pm$ 0.0002047 \\

FieldFormer
& \textbf{0.09888 $\pm$ 0.0007416}
& \textbf{0.07876 $\pm$ 0.0005815}
& \textbf{0.1154 $\pm$ 0.00032}
& 0.09216 $\pm$ 0.0002176
& \textbf{0.04644 $\pm$ 0.000254}
& \textbf{0.03703 $\pm$ 0.0001957} \\

FieldFormer-Gamma Field
& 0.225 $\pm$ 0.001009
& 0.1741 $\pm$ 0.0006781
& \textbf{0.1154 $\pm$ 0.0002949}
& \textbf{0.09212 $\pm$ 0.0001884}
& 0.04696 $\pm$ 0.0002675
& 0.03738 $\pm$ 0.000196 \\

FieldFormer-MLP
& 0.09908 $\pm$ 0.000686
& 0.07895 $\pm$ 0.0005192
& 0.1157 $\pm$ 0.0003294
& 0.09245 $\pm$ 0.0002116
& 0.04713 $\pm$ 0.0002381
& 0.03761 $\pm$ 0.000187 \\

\hline
\end{tabular}%
}
\end{table}

Table~\ref{tab:ablation_synthetic_global} reports global field reconstruction results on the synthetic datasets. Unlike sensor-space imputation, the best variant changes across datasets. FieldFormer-MLP performs best on Heat, FieldFormer and FieldFormer-Gamma Field are strongest on different Pollution metrics, and FieldFormer-Gamma Field gives the best global reconstruction on SWE among FieldFormer variants. This supports the broader observation that global reconstruction is more prior-dependent than sensor-space imputation: architectural changes that help away from sensors do not necessarily improve sensor-space prediction.

\begin{table}[htbp]
\centering
\caption{Ablation study on synthetic PDE benchmarks: global field reconstruction. Dense reference fields are available for these datasets. Lower is better.}
\label{tab:ablation_synthetic_global}
\resizebox{\textwidth}{!}{%
\begin{tabular}{lcccccc}
\hline
Model 
& Heat RMSE & Heat MAE
& Pollution RMSE & Pollution MAE
& SWE RMSE & SWE MAE \\
\hline

FieldFormer-PINN
& 0.1641 $\pm$ 0.001102
& 0.1189 $\pm$ 0.0007892
& 0.4582 $\pm$ 0.04409
& 0.1386 $\pm$ 0.00625
& 0.2374 $\pm$ 0.00019
& 0.1717 $\pm$ 0.000157 \\

FieldFormer
& 0.1818 $\pm$ 0.001179
& 0.1313 $\pm$ 0.0008449
& \textbf{0.4359 $\pm$ 0.04525}
& 0.0831 $\pm$ 0.006115
& 0.2015 $\pm$ 0.0002546
& 0.1465 $\pm$ 0.0002115 \\

FieldFormer-Gamma Field
& 0.2193 $\pm$ 0.001191
& 0.1641 $\pm$ 0.0008903
& 0.4366 $\pm$ 0.04515
& \textbf{0.08203 $\pm$ 0.006125}
& \textbf{0.1874 $\pm$ 0.0001986}
& \textbf{0.1362 $\pm$ 0.0001916} \\

FieldFormer-MLP
& \textbf{0.1594 $\pm$ 0.0009926}
& \textbf{0.1144 $\pm$ 0.0007596}
& 0.4645 $\pm$ 0.04324
& 0.1455 $\pm$ 0.006278
& 0.219 $\pm$ 0.000232
& 0.1577 $\pm$ 0.0001574 \\

\hline
\end{tabular}%
}
\end{table}

\paragraph{Real-world sparse sensor benchmarks.}
For real-world datasets, we compare FieldFormer against FieldFormer-Gamma Field and FieldFormer-MLP, while excluding PINN variants because no reliable real-world physics residual is available. Table~\ref{tab:ablation_real_sparse} shows sensor-space imputation results on atmospheric and pollution monitoring datasets. FieldFormer-MLP tests whether sparse local neighborhood construction alone is sufficient when the local transformer is replaced by an MLP aggregator. Gamma Field improves several wind and pollution metrics, but does not uniformly improve all atmospheric variables. In particular, FieldFormer remains better for relative humidity and some mean absolute error metrics. This suggests that coordinate-dependent neighborhood geometry can help when local scales vary spatially, but also introduces additional flexibility that may not always be beneficial under sparse real-world sensing.

\begin{table}[htbp]
\centering
\caption{Ablation study on real-world benchmarks: sensor-space imputation. Results are reported as mean $\pm$ std over bootstrap samples. Lower is better.}
\label{tab:ablation_real_sparse}
\resizebox{0.95\textwidth}{!}{%
\begin{tabular}{llcccccc}
\hline
Dataset & Variable 
& FieldFormer RMSE & FieldFormer MAE
& Gamma Field RMSE & Gamma Field MAE
& FieldFormer-MLP RMSE & FieldFormer-MLP MAE \\
\hline

Atmospheric & AT
& 0.9966 $\pm$ 0.02692
& \textbf{0.5198 $\pm$ 0.004008}
& \textbf{0.9935 $\pm$ 0.02355}
& 0.5232 $\pm$ 0.004039
& 1.025 $\pm$ 0.02321
& 0.5712 $\pm$ 0.004604 \\

Atmospheric & RH
& \textbf{4.333 $\pm$ 0.0544}
& \textbf{2.171 $\pm$ 0.01099}
& 4.498 $\pm$ 0.07141
& 2.18 $\pm$ 0.01716
& 4.425 $\pm$ 0.04953
& 2.267 $\pm$ 0.01041 \\

Atmospheric & $U_x$
& 0.6807 $\pm$ 0.0134
& 0.3213 $\pm$ 0.002875
& \textbf{0.6547 $\pm$ 0.009111}
& \textbf{0.3189 $\pm$ 0.002666}
& 0.6888 $\pm$ 0.01342
& 0.3436 $\pm$ 0.003317 \\

Atmospheric & $V_y$
& 0.7011 $\pm$ 0.01045
& 0.3468 $\pm$ 0.002512
& \textbf{0.6997 $\pm$ 0.009647}
& \textbf{0.3457 $\pm$ 0.002271}
& 0.7388 $\pm$ 0.009629
& 0.3717 $\pm$ 0.00248 \\

Pollution & PM$_{10}$
& 39.25 $\pm$ 0.8496
& 20.28 $\pm$ 0.1172
& \textbf{38.83 $\pm$ 0.9117}
& \textbf{19.91 $\pm$ 0.1438}
& 40.23 $\pm$ 0.7156
& 21.23 $\pm$ 0.1219 \\

Pollution & PM$_{2.5}$
& 22.06 $\pm$ 0.664
& 11.23 $\pm$ 0.0691
& \textbf{21.99 $\pm$ 0.72}
& \textbf{11.01 $\pm$ 0.07286}
& 23.24 $\pm$ 0.6644
& 11.65 $\pm$ 0.06181 \\

\hline
\end{tabular}%
}
\end{table}


Overall, the ablations show that FieldFormer’s core locality-aware architecture is already strong for sensor-space imputation. Physics regularization does not consistently improve sparse prediction and can degrade performance when the residual objective conflicts with noisy or underconstrained observations. Replacing the transformer with an MLP preserves some benefits of adaptive neighborhood construction and remains competitive on several real-world variables, but it underperforms the transformer-based variants on most metrics. This suggests that attention-based relational aggregation improves reliability in heterogeneous real-world sensor data. Gamma Field improves selected global and heterogeneous real-world metrics, especially for some pollution and wind variables, but its gains are not uniform, suggesting that coordinate-dependent geometry is useful but should be evaluated on a case-by-case basis.
